\section{Introduction}

When we evaluated an open-weight language model on the ANLI R1 benchmark --- a dataset specifically constructed so that human annotators could fool state-of-the-art NLP models --- we found that the model's Expected Calibration Error \textit{decreased} relative to its performance on clean MultiNLI. The adversarial benchmark, designed to expose model failures, made the model appear \textit{better calibrated}. This result was not an error. It reveals a fundamental property of how adversarial construction methods interact with model confidence: in a model-in-the-loop adversarial benchmark, the examples are selected precisely because the model fails on them --- and a model that fails with appropriately low confidence is, by definition, well-calibrated.

This observation motivates the central question of our work: \textbf{does adversarial perturbation of NLP benchmarks actually degrade LLM calibration, and if so, under what conditions?}

\subsection{The Problem}

Calibration --- the alignment between predicted confidence and observed accuracy --- is a critical reliability property for deployed language models. A model that assigns 90\% confidence to its predictions should be correct 90\% of the time; deviations from this correspondence make model confidence unreliable as a safety signal. The field has developed robust measurement tools, most notably Expected Calibration Error (ECE) \citep{guo2017calibration}, and has documented that modern neural networks are systematically overconfident under distribution shift \citep{minderer2021revisiting}.

At the same time, adversarial NLP benchmarks --- AdvGLUE \citep{wang2021adversarial} and ANLI \citep{nie2020adversarial} --- have become standard for evaluating LLM robustness under distribution shift, typically reporting accuracy degradation. But accuracy degradation alone does not characterize the reliability of model confidence under adversarial conditions. A model that becomes \textit{less} accurate but \textit{proportionally less confident} may be better calibrated than one that maintains confidence while failing more often.

The deeper problem is that adversarial benchmarks are not a single category. AdvGLUE uses \textit{human adversaries} who craft perturbations to fool existing models; ANLI uses a \textit{model-in-the-loop} construction where examples are selected by human annotators specifically when they cause model errors. These construction methods produce qualitatively different distributions of model behavior --- and consequently, different calibration effects. Yet prior calibration studies treat adversarial perturbation as a monolithic distribution shift, analogous to the vision-domain corruptions studied by \citet{minderer2021revisiting}. Whether this analogy holds in NLP, and whether RLHF alignment moderates the effect, has not been tested.

\textbf{The gap:} No published work measures logit-based ECE on adversarial NLP benchmark splits (AdvGLUE, ANLI) for open-weight LLMs, or examines how adversarial construction method and RLHF alignment jointly determine calibration outcomes under adversarial conditions.

\subsection{Key Insight}

Our central finding is that adversarial calibration degradation in LLMs is \textbf{task-type and construction-method conditional}, not universal. Human-adversarial NLI examples (AdvGLUE MNLI) produce reliable calibration degradation: $\Delta$ECE $= +0.071$, a 7.1 percentage-point increase in Expected Calibration Error relative to clean MultiNLI. Model-in-the-loop adversarial NLI examples (ANLI R1/R2) produce calibration \textit{improvement}: $\Delta$ECE $= -0.041$ for R1. Binary classification tasks (QQP, SST-2) show reversed $\Delta$ECE regardless of adversarial method. And RLHF alignment conditionally moderates these effects --- helping for model-in-loop adversarial (100\% moderation rate across ANLI rounds) but exacerbating calibration degradation for static human-adversarial benchmarks ($\Delta\Delta$ECE $= -0.026$, chat worse than base on AdvGLUE).

The resolution to the opening paradox follows directly: in model-in-the-loop adversarial benchmarks, examples are selected to cause errors; a model that fails with appropriately moderate confidence on such examples is demonstrating calibration, not miscalibration. Human-adversarial examples, by contrast, exploit surface features to trigger high confidence on wrong answers --- the classic miscalibration mechanism.

\subsection{Contributions}

Our investigation of adversarial LLM calibration yields three contributions that build from existence to mechanism to conditionality:

\textbf{(1) First measurement of logit-based ECE on adversarial NLP benchmarks.} We measure ECE directly on AdvGLUE and ANLI benchmark splits for Llama-2-7b-hf, establishing $\Delta$ECE as a measurable, meaningful quantity for adversarial calibration stress testing. This fills the gap between adversarial robustness evaluation (accuracy-only) and calibration measurement (clean-data-only). Our JSONL logit caching strategy enables multiple downstream calibration analyses (label preservation, confidence-accuracy decoupling, per-cell $\Delta$ECE, RLHF comparison) from a single inference pass, making multi-cell adversarial calibration analysis computationally feasible.

\textbf{(2) Evidence that calibration degradation is task-type and construction-method conditional.} Human-adversarial NLI examples (AdvGLUE MNLI: $\Delta$ECE $= +0.071$) produce calibration degradation; model-in-the-loop adversarial examples (ANLI R1/R2: $\Delta$ECE $< 0$) produce calibration improvement; and the ANLI difficulty gradient (R1 $\rightarrow$ R2 $\rightarrow$ R3) shows a smooth transition from improvement to degradation as adversarial difficulty increases. Binary classification tasks consistently show reversed $\Delta$ECE. This nuances the vision-domain finding that distribution shift uniformly degrades calibration.

\textbf{(3) Evidence that RLHF alignment conditionally moderates adversarial calibration degradation.} We compare Llama-2-7b-hf (base) and Llama-2-7b-chat across ANLI and AdvGLUE conditions. RLHF moderation is confirmed for model-in-loop adversarial (100\% of ANLI rounds show $\Delta\Delta$ECE $> 0.01$ favoring chat), but reversed for static human-adversarial benchmarks (AdvGLUE: chat shows larger $\Delta$ECE than base by 2.6 percentage points). This is the first documented benchmark-type $\times$ RLHF calibration interaction.

Together, these contributions motivate adversarial calibration measurement as a distinct research problem from adversarial accuracy robustness, and provide a conditional framework (construction method $\times$ task type $\times$ RLHF alignment) for understanding when and how adversarial perturbation may degrade LLM reliability signals --- within the scope of this single-model pilot study.

We organize the paper as follows. Section~\ref{sec:related} discusses calibration measurement, adversarial NLP benchmarks, and RLHF calibration --- three bodies of work whose intersection we occupy. Section~\ref{sec:methodology} describes our measurement methodology. Section~\ref{sec:experiments} presents our experimental design. Section~\ref{sec:results} reports results across the four research questions. Section~\ref{sec:discussion} interprets findings and discusses limitations. Section~\ref{sec:conclusion} concludes.
